% ============================================================
% DRAFT — proposed additions to chapters/introduction.tex, aimed
% at the specific examiner questions you flagged from your own
% voice notes: "why didn't you just use an 8/10GB Jetson on the
% rover" and "why did you make the laptop/desktop smart and the
% rover dumb instead of distributing intelligence onboard".
%
% The intro already motivates local-only processing (privacy,
% cost, no LiDAR/IMU) via the "private home" framing in
% Section~\ref{sec:background}. What it does NOT yet do is name
% the rejected alternative (an onboard embedded accelerator) and
% say explicitly why the split was chosen instead. That's what
% these blocks add. Nothing here has been written into the actual
% chapter file — read it next to introduction.tex and merge
% whatever you're happy with.
%
% J1. New paragraph at the end of Section~\ref{sec:background},
%     right before "\mysection{Problem Statement...}", naming the
%     rejected onboard-Jetson alternative and giving the four
%     reasons it was rejected: already-owned-hardware premise,
%     compute headroom, benchmark validity (RQ1 needs the model to
%     be the only variable), and generality (rover as instrument,
%     not as the thing under test).
% J2. Small wording change to RQ1 so the host/rover split is
%     baked into the question itself, not just the surrounding
%     prose — so it can't be read as an oversight.
% J3. New paragraph in Section~\ref{sec:scope-and-limitations},
%     stating plainly that the rover's own onboard computer
%     (the Pi) is treated as instrumentation, not as part of the
%     system under test.
% J4. A couple of phrases in J1 are flagged [[VERIFY]] — specific
%     cost/TOPS figures for Jetson Orin Nano/NX 8-16GB modules
%     at time of writing. I didn't want to assert numbers I
%     can't verify from here; the qualitative argument holds
%     without them, but citing a real spec sheet or current price
%     would make it land harder in the oral.
% ============================================================


% ---------- J1. New paragraph, end of Section~\ref{sec:background} ----------
% CURRENT, this is the last paragraph of sec:background (currently ends the
% section right before "\mysection{Problem Statement and Research Questions}"):
%   ...and the robot has to rely on its camera to know whether its plan is working. Where privacy is not
%   a concern and a reliable internet connection is available, a cloud-hosted model remains the
%   better choice, since it is larger and more capable than any model that fits on consumer
%   hardware.
% INSERT the following as a NEW paragraph immediately after that one (still
% inside sec:background, still before the Problem Statement section header):

An alternative to running the models on a separate host machine would be to mount a dedicated
embedded AI accelerator, such as an NVIDIA Jetson Orin Nano or Orin NX, directly on the rover and
run the entire pipeline on board. This project deliberately does not do that, for four reasons.
First, it would undercut the premise that motivates local processing in the first place: the home
robot framing above assumes the household already owns the compute the system needs, in the form
of a desktop or laptop; requiring the purchase of a separate embedded accelerator module
[[VERIFY: current price, e.g.\ USD~X--Y]] on top of the rover itself reintroduces exactly the kind
of extra, specialised hardware cost this project is trying to avoid. Second, even the larger
embedded modules in this class have a fraction of the memory and raw throughput of a consumer
desktop GPU [[VERIFY: e.g.\ Jetson Orin NX 16GB TOPS vs.\ the GTX~Titan~X's compute, or simply
cite both spec sheets]], so moving the same models onto the rover itself would force either a
smaller, weaker model, or a slower response time, or both, confounding what RQ1 and RQ2 are
actually trying to measure: the capability limits of the models, not the engineering limits of an
embedded board. Third, keeping inference off the rover, on a host machine that can run the same
planner locally or swap to a cloud API, is what makes the local-versus-cloud comparison in RQ1 an
apples-to-apples comparison of model capability alone; if the rover itself hosted the model,
switching the benchmark to a cloud API would also change the hardware story, confounding the
comparison. Fourth, it keeps the result general: the rover's own on-board computer, a Raspberry
Pi, only has to stream camera frames and relay drive and sensor commands over the local network,
so the same rover would work unchanged with a different host machine entirely, a laptop today, a
household's existing desktop tomorrow, rather than being tied to whatever accelerator happens to
be bolted onto this specific rover.


% ---------- J2. RQ1 wording — bake the host/rover split into the question itself ----------
% CURRENT (inside the \begin{enumerate} block under
% "\mysection{Problem Statement and Research Questions}"):
%   \item[\textbf{RQ1}] Can locally hosted models, running on the consumer hardware defined
%       above, turn a vague spoken command into a plan that the rover can execute, and which
%       levels of command in Table~\ref{tab:command-levels} can they handle reliably?
% REPLACE WITH:

\item[\textbf{RQ1}] Can locally hosted models, running on the consumer host machine defined
    above rather than on the rover's own on-board computer, turn a vague spoken command into a
    plan that the rover can execute, and which levels of command in Table~\ref{tab:command-levels}
    can they handle reliably?


% ---------- J3. New paragraph in Section~\ref{sec:scope-and-limitations} ----------
% CURRENT, the paragraph that currently reads:
%   The rover is deliberately limited to a camera, an ultrasonic distance sensor and a line sensor,
%   with no wheel encoders, \ac{imu} or \acs{lidar}. This reflects the low-cost home robot described
%   in Section~\ref{sec:background}, and forces the pipeline to rely on the camera to determine
%   whether a plan is succeeding, which is the capability this project sets out to test. As a result,
%   movement is carried out without precise measurement of distance or heading, and commands that
%   require precise distances beyond what the ultrasonic sensor can measure are not supported.
% INSERT the following as a NEW paragraph immediately after that one:

In the same spirit, the rover's own on-board computer is treated as instrumentation rather than
as part of the system under test. It streams camera frames to the host machine and relays the
drive and sensor commands it receives back, but performs no part of the planning, perception or
clarification pipeline itself. The language and vision models that answer RQ1 and RQ2 all run on
the host machine defined in Section~\ref{sec:background}, so that both research questions probe
the capability of those models on fixed, known hardware, rather than being confounded by whatever
embedded computer happens to be fitted to the rover.

